% !TeX program = lualatex
\makeatletter
\def\input@path{{00/plantilla/formato-base/}{04/}}
\makeatother
\documentclass[carta,firma]{lafrox}
\usepackage{etoolbox}
\usepackage{placeins}
\captionsetup{hypcap=false}
\renewcommand{\LFXisotipoportada}{00/plantilla/formato-base/logo/lafrox_isotipo_naranja.pdf}
\renewcommand{\LFXisotipopagina}{00/plantilla/formato-base/logo/lafrox_isotipo_negro.pdf}
% Comando de procedencia de figuras usado por 4.2 y 4.3; la plantilla base no lo incluye.
\newcommand{\fuentefigura}[1][elaboración propia]{%
  \par\vspace{2pt}%
  {\footnotesize\color{lafroxGris}\raggedright Fuente: #1.\par}}
\datoslafrox{
  documento = {Propuesta Técnica},
  subtitulo = {Subdocumento 4 --- Arquitectura de la solución},
  alcance = {Apartados 4.1, 4.2 y 4.3},
  formulario = {Formulario T-7},
  fecha = {05 de octubre de 2026},
  version = {},
  nomenclatura = {LAFROX-Subdocumento4.pdf},
}
\hypersetup{pdftitle={Subdocumento 4 -- Arquitectura de la solución}}
\renewcommand{\LFXestadodoc}{Entrega 2}
\graphicspath{{04/}}
% Numeración corrida de tablas y figuras en todo el subdocumento.
\counterwithout{table}{chapter}
\counterwithout{figure}{chapter}

\begin{document}
\portadalafrox
\preliminareslafrox
\setcounter{chapter}{3}
% Ensamblador del Subdocumento 4: 4.1, 4.2 y 4.3, seguidos de Referencias y Declaración de uso de IA.
% La numeración de tablas y figuras es corrida en todo el subdocumento (main.tex).
\begingroup
\import{04/partes/4.1_logica/}{contenido.tex}
\endgroup
\clearpage
\begingroup
\import{04/partes/4.2_fisica/}{contenido.tex}
\endgroup
\clearpage
\import{04/partes/cierre/}{referencias.tex}
\import{04/partes/cierre/}{declaracion_de_uso_de_ia.tex}

\end{document}
